\documentclass[10pt,twocolumn]{article}
\usepackage[scaled]{helvet}
\renewcommand\familydefault{\sfdefault}
\usepackage[T1]{fontenc}
\usepackage[margin=0.7in]{geometry}
\usepackage{titlesec}
\usepackage{enumitem}
\usepackage{parskip}
\setlength{\columnsep}{0.24in}
\setlist[itemize]{leftmargin=1.2em, topsep=0.2em, itemsep=0.18em}
\titleformat{\section}{\large\bfseries}{\thesection}{1em}{}

\begin{document}
\title{\vspace{-1.6cm}AWS S3 (Simple Storage Service)}
\author{AWS ML Study Notes}
\date{}
\maketitle
\small

\section*{Overview}
S3 is AWS object storage. It stores data as objects inside buckets and is designed for very high durability, massive scale, and broad integration with the rest of AWS.

In ML systems, S3 is usually the default data lake and artifact store. Raw data, curated datasets, model files, pipeline outputs, logs, and static assets often live there.

\section*{Core Concepts}
\begin{itemize}
\item Buckets are top level containers and objects are individual files or blobs addressed by key. Prefixes are naming conventions that simulate folders and support partitioned layouts.
\item Important features include versioning, lifecycle rules, server side encryption, bucket policies, replication, event notifications, and multiple storage classes.
\item S3 provides strong read after write and list consistency, which simplifies many ingestion and pipeline patterns compared with older eventual consistency assumptions.
\end{itemize}

\section*{How It Works}
\begin{itemize}
\item Applications upload or download objects through the S3 API, SDKs, CLI, or signed URLs. IAM and bucket policies decide who can access which paths.
\item Data engineering jobs commonly read from raw prefixes, transform the data, and write curated outputs in columnar formats such as Parquet for analytics and downstream ML use.
\item Many services integrate directly with S3, including Athena, Glue, SageMaker, Lambda event notifications, and CloudTrail log delivery.
\end{itemize}

\section*{AI and ML Relevance}
\begin{itemize}
\item Training datasets are often staged in S3, and model artifacts are written back there after the training job completes.
\item S3 can trigger downstream workflows when new data lands, which is useful for automated feature computation, validation, batch inference, or retraining pipelines.
\end{itemize}

\section*{Operational Notes}
\begin{itemize}
\item Design the bucket key structure deliberately. Bad partitioning and inconsistent naming make analytics, retention, and lifecycle management painful later.
\item Encrypt sensitive data, control public access at both the bucket and account level, and use lifecycle rules to move cold data to cheaper storage classes.
\item Do not treat S3 like a filesystem with heavy small file chatter when a compact columnar layout would be more efficient and cheaper.
\end{itemize}

\section*{What To Remember}
\begin{itemize}
\item S3 is object storage, not block storage and not a database.
\item Bucket policy plus IAM policy plus encryption controls together define access posture.
\item For ML, think of S3 as the durable backbone for data and artifacts rather than just a place to drop files.
\end{itemize}

\end{document}
